\chapter*{Streszczenie}

Celem niniejszej pracy magisterskiej jest eksperymentalna ocena wpływu mechanizmu adaptacyjnej ewolucji kontekstu agentowego na skuteczność autoremediacji aplikacji mikrousługowych w środowisku Kubernetes. W pracy zaprojektowano hierarchiczny system wieloagentowy oparty na paradygmacie \textit{agent-as-a-tool}, w którym \textit{Incident Commander} koordynuje pracę trzech wyspecjalizowanych deputatów, wyposażony w mechanizm samouczenia oparty na algorytmie \textit{Agentic Context Engineering} (ACE), rozszerzonym o hierarchiczną propagację refleksji wzdłuż drzewa zadań oraz mechanizm ukrywania wag skuteczności reguł przed agentem wykonawczym.

Eksperymenty przeprowadzono na izolowanym środowisku z referencyjną aplikacją mikrousługową \textit{Robot-shop} i autorskim komponentem \textit{Saboteur Agent}, generującym deterministyczną sekwencję 98 incydentów reprezentujących cztery klasy usterek. Badania objęły dwa modele językowe - \textit{GPT-5-nano} i \textit{Minimax-2.5} - przy identycznej sekwencji epizodów, z automatyczną oceną trzech binarnych wskaźników sukcesu.

Wyniki potwierdzają główną tezę pracy: \textit{GPT-5-nano} osiąga poprawę skuteczności remediacji o $+18{,}4$~pp ($+120\%$ względnie), \textit{Minimax-2.5} o $+5{,}1$~pp ($+13\%$ względnie). Modele wykształcają odmienne strategie organizacyjne - GPT-5-nano uczy się koordynacji agentów, Minimax-2.5 uczy się zamykania cyklu remediacji przez wdrożenie. Zaobserwowano zjawisko emergentne: średnia maksymalna głębokość drzewa zadań rośnie o $+26{,}1\%$ bez jawnego programowania tej strategii.

\vspace{20pt}
\noindent Słowa kluczowe: \textit{Site Reliability Engineering}, autoremediacja, systemy wieloagentowe, duże modele językowe, \textit{Agentic Context Engineering}, samouczenie agentów, mikrousługi, Kubernetes, drzewo zadań
